\section{Error and Error Rate}

\subsection{Error}

The {\bf error} $\ve{e}(t)$ between two signals 
$\va(t)$ and $\vb(t)$ is 
defined as the difference between the two signals, 
\be
\ve{e}(t) \defe \abs{\va(t) - \vb(t)}.
\ee
Note that $\ve{e} \geq \ve{0} \;\; \forall \;\; t$.
If $\vb$ exactly follows $\va$ at every time step, then the
error is zero, as expected.  

\begin{Hexample}
{
Let vectors $\va$ and $\vb$ be defined as follows,
\begin{align}
 \va & = [\; 20, 20, 20, 20 \;] \\
 \vb & = [\; 18, 18, 18, 18 \;] 
\end{align}
Then the error $\ve{e}$ is found to be,
\be
 \ve{e} = [\; 2, 2, 2, 2 \;]
\ee
In this example, assuming $\va$ is the ground truth, $\vb$ 
underestimates $\va$ by $2$ at every sample.  This is an 
example of {\bf systemic error}, where there is a consistent
offset between the two signals, and an offset adjustment (the
offset value of $2$ in this case) would eliminate all error 
between the two signals.
}
\end{Hexample}

\subsection{Error Rate}

Next, we would like to have a scalar value to describe the 
error over the entire sample, which is measured at each sample
point and has dimension $n$. 

We defined the {\bf accumulated error}, 
$\sum \ve{e}$, as the sum of the errors
of each of the measurement,
\be
\sum \ve{e} \defe \sum_{i=0}^{n-1} \abs{e_i} = 
\abs{a_i - b_i}.
\ee
A serious problem with accumulated error is that the error 
increases the more times the signals are measured. 

\begin{Hexample}
{
\begin{itemize}
\item In the preceding example, 
if the signals were measured at only 
one sample, the accumulated error would be $2$.  
\item Similarly, 
if two samples were taken, the accumulated error would be $4$.
\item For three samples, the accumulated error would be $6$.  
\item For 
the example as stated above, with four samples, the accumulated
error is $8$.  
\end{itemize}
This pattern shows the limitation of accumulated error: that
it is dependent on sample size.
}
\end{Hexample}
